\page
\setupbackgrounds[page][background=appendixfurniture]
\switchtobodyfont[11pt]
\setupinterlinespace[line=15pt]
\setupalign[flushleft,nothyphenated,verytolerant,stretch]
{\switchtobodyfont[20pt]\color[accent]{O vocabulário de Persas}}\par
\blank[5mm]
{\bf Levantamento preliminar}\par
\blank[2mm]
Este levantamento reúne 1.639 lemas em 4.755 ocorrências analisadas, sem nomes próprios.
\par\blank[3mm]
\starttabulate[|p(56mm)|r|r|]
\NC \bf Categoria \NC \bf Lemas \NC \bf Ocorrências \NC\NR
\HL
\NC básico DCC\NC 333 (20,3\%) \NC 2604 (54,8\%) \NC\NR
\NC recorrente na tragédia\NC 885 (54,0\%) \NC 3815 (80,2\%) \NC\NR
\NC preferência por Ésquilo\NC 99 (6,0\%) \NC 208 (4,4\%) \NC\NR
\NC mais frequente fora da tragédia\NC 208 (12,7\%) \NC 570 (12,0\%) \NC\NR
\stoptabulate
\blank[3mm]
As categorias se sobrepõem: seus percentuais não devem ser somados. Lema é a entrada do dicionário; ocorrência é cada aparição de uma forma no texto.\par
\blank[3mm]
{\bf Como ler as categorias}\par
\blank[2mm]
{\bf Básico DCC.} Pertencimento à lista de 500 palavras do Dickinson College Commentaries; os 24 numerais suplementares foram excluídos. É uma referência pedagógica, não uma medida absoluta de facilidade.\par
\blank[2mm]
{\bf Recorrente na tragédia.} Presença em pelo menos três outras peças e em dois autores da amostra. Esta recorrência aproxima um núcleo trágico; não prova que uma palavra seja fácil para o leitor.\par
\blank[2mm]
{\bf Preferência por Ésquilo.} Pelo menos três ocorrências em duas outras peças atribuídas a Ésquilo, com frequência relativa quatro vezes maior que na amostra de Sófocles e Eurípides.\par
\blank[2mm]
{\bf Mais frequente fora da tragédia.} Pelo menos cinco ocorrências em um contexto comparado e frequência relativa quatro vezes maior que na amostra trágica. Indica distribuição lexical, não a origem ou o sentido da palavra nesta passagem.\par
\page
{\switchtobodyfont[18pt]\color[accent]{Corpus e limites}}\par
\blank[4mm]
{\bf O texto analisado}\par
\blank[2mm]
Das 5217 unidades gráficas do texto de Sidgwick, 4755 receberam lema comum, 248 foram identificadas como nomes próprios e 214 ficaram sem classificação. A cobertura de identificação é 95,9\%.\par
\blank[2mm]
As formas foram confrontadas com o Perseus Treebank de Persas, baseado em outra edição. Exigimos uma análise lexical única na posição mais próxima, a no máximo três versos. Divergências, ambiguidades e fragmentos ficam fora dos denominadores. Crases e elisões seguem a segmentação gráfica e os lemas disponíveis no treebank.\par
\blank[3mm]
{\bf Amostra de comparação}\par
\blank[2mm]
Persas foi excluída da amostra de referência. Usamos seis outras peças atribuídas a Ésquilo (incluindo Prometeu acorrentado, de autoria discutida), cinco de Sófocles e Medeia, de Eurípides. A amostra de Eurípides é pequena e a comparação entre autores permanece provisória.\par
\blank[2mm]
Fora da tragédia: Ilíada e Odisseia; Acarnenses e Tesmoforiantes; o trecho anotado da Poliorcética de Eneias Tático; De automatis de Herão; trechos e obras de Longo, Luciano e o discurso 24 de Lísias. São amostras de extensão e época diferentes.\par
\blank[2mm]
As frequências são divididas pelo total de ocorrências comuns de cada corpus. Nos testes de preferência, acrescentamos 0,5 ao número de ocorrências para estabilizar razões quando há zeros. Não confundimos ausência na amostra com exclusividade ou criação de Ésquilo.\par
\blank[3mm]
Para as margens, adotamos um recorte editorial mais restrito: 175 lemas fundamentais (os primeiros 150 da lista de frequência DCC e 25 acréscimos). O sentido básico é dispensado; formas difíceis podem conservar análise gramatical. Esse recorte não equivale à categoria de 500 palavras da tabela.\par
\blank[2mm]
{\bf Fontes e reprodução}\par
\blank[2mm]
Dickinson College Commentaries: Greek Core Vocabulary (CC BY-SA 3.0). Perseus Dependency Treebank 2.1 (CC BY-SA 3.0 US). Pedalion Trees: Toon Van Hal, Alek Keersmaekers e colaboradores (CC BY-SA 4.0). Dados consultados em 10 de outubro de 2026.\par
\blank[2mm]
O projeto conserva a lista de lemas, as contagens por corpus, a classificação de cada unidade e o programa de cálculo. Este apêndice será atualizado depois da revisão filológica das análises e da ampliação dos corpora.\par
